\chapter{بازسازی و پیاده‌سازی روش مرجع}
\label{ch:reference}

\section{مقدمه}

در این فصل، روش مرجع ارائه‌شده در \cite{chen2024dsp} برای حل مسئله مجموعه احاطه‌گر با استفاده از شبکه عصبی گرافی و یادگیری تقویتی دوگانه بازسازی و پیاده‌سازی می‌شود. هدف، اطمینان از درک صحیح چارچوب و ایجاد پایه‌ای برای گسترش آن به مسئله مجموعه احاطه‌گر کامل است.

\section{مدلسازی مسئله به‌صورت فرآیند تصمیم‌گیری مارکوف}

روش مرجع، مسئله مجموعه احاطه‌گر را به‌صورت یک فرآیند تصمیم‌گیری ترتیبی مدلسازی می‌کند. در هر گام، عامل یکی از رأس‌های مجاز را انتخاب می‌کند تا به مجموعه $S$ افزوده شود. حالت محیط با بردار ویژگی هر رأس و ساختار گراف توصیف می‌شود \cite{chen2024dsp}.

\section{معماری شبکه عصبی گرافی}

شبکه عصبی به‌کاررفته در این روش از نوع MPNN است. نمایش اولیه هر رأس با بردار ویژگی‌ای شامل اطلاعاتی نظیر درجه رأس و وضعیت پوشش آن ساخته می‌شود. سپس با عبور از چند لایه پیام‌رسانی، نمایش‌های غنی‌تری تولید شده و در نهایت برای پیش‌بینی مقادیر $Q$ هر عمل مجاز به‌کار می‌روند \cite{gilmer2017mpnn, chen2024dsp}.

\section{الگوریتم یادگیری تقویتی دوگانه}

برای آموزش عامل از الگوریتم $Double\ DQN$ استفاده می‌شود. حافظه تجربه\footnote{Replay Buffer} برای ذخیره انتقال‌ها و نمونه‌برداری تصادفی جهت شکستن همبستگی، به‌کار می‌رود. همچنین هدف‌گیری آرام\footnote{Soft Target Update} برای افزایش پایداری یادگیری اعمال می‌گردد \cite{vanhasselt2016ddqn, mnih2015human}.

\section{طراحی فضای عمل و ماسک عمل}

در روش مرجع، در هر حالت تنها رأس‌های «مجاز» می‌توانند انتخاب شوند. این محدودسازی باعث کاهش فضای جستجو و تسریع یادگیری می‌گردد. انتخاب رأس‌های نامجاز طبق قواعد تعریف‌شده در محیط، از طریق ماسک عمل\footnote{Action Masking} غیرفعال می‌شود \cite{chen2024dsp}.

\section{پیاده‌سازی و بازتولید نتایج}

با پیاده‌سازی دقیق چارچوب فوق، تلاش شد نتایج گزارش‌شده در \cite{chen2024dsp} روی خانواده‌های گراف استاندارد بازتولید گردد. این بازتولید، پایه لازم برای تحلیل صحت پیاده‌سازی و مقایسه بعدی با نسخه گسترش‌یافته به مسئله مجموعه احاطه‌گر کامل را فراهم می‌کند.

\section{خلاصه}

در این فصل، روش مرجع مبتنی بر $GNN + DDQN$ برای مسئله مجموعه احاطه‌گر بازسازی شد. مفاهیم کلیدی شامل مدلسازی ترتیبی، MPNN، $DDQN$، حافظه تجربه و ماسک عمل بیان گردید. این پیاده‌سازی، نقطه شروع گسترش به مسئله مجموعه احاطه‌گر کامل در فصل بعدی خواهد بود.
